\chapter{Họp Phụ Huynh Bão Táp}
\markboth{Họp Phụ Huynh Bão Táp}{Họp Phụ Huynh Bão Táp}

\section{Trên Lớp Một Kiểu, Trước Phụ Huynh Một Kiểu}
\begin{truyen}{Trên Lớp Một Kiểu, Trước Phụ Huynh Một Kiểu}{One Face in Class, Another Before Parents}
\chuhoa{H}{ọc sinh:} ``Em sợ họp phụ huynh hơn cả kiểm tra, vì sau đó kiểu gì nhà em cũng thành chiến trường.''
\textit{(Student: ``I fear parent meetings more than exams, because afterwards my house always turns into a battlefield.'')}

\teacherpair{Nỗi sợ đó có thật ở nhiều học sinh. Họp phụ huynh lẽ ra để nối nhà trường với gia đình. Nhưng nếu nó chỉ biến thành buổi tổng kết lỗi lầm và so điểm, thì sau buổi họp, kiến thức chưa chắc tăng mà tổn thương trong nhà lại tăng rất nhanh.}{``That fear is real for many students. Parent meetings should connect school and family. But if they become mere sessions of listing mistakes and comparing scores, then after the meeting knowledge may not rise while emotional damage at home rises quickly.''}
\end{truyen}

\section{Phụ Huynh Bênh Con Mù Quáng}
\begin{truyen}{Phụ Huynh Bênh Con Mù Quáng}{Parents Who Defend Blindly}
\chuhoa{H}{ọc sinh:} ``Có mấy bạn cứ yên tâm sai vì biết ba mẹ sẽ đứng ra cãi giúp.''
\textit{(Student: ``Some students feel safe doing wrong because they know their parents will argue for them anyway.'')}

\teacherpair{Đó cũng là một kiểu làm hỏng con rất tinh vi. Bảo vệ con khỏi sự bất công là cần. Nhưng bảo vệ con khỏi hậu quả của lỗi sai mà chính nó gây ra thì đang cắt mất cơ hội trưởng thành. Một đứa trẻ luôn có người cãi thay rất dễ lớn thành người không biết chịu trách nhiệm.}{``That is another subtle way of damaging a child. Protecting children from injustice is necessary. But protecting them from the consequences of their own mistakes cuts away their chance to grow. A child who is always defended grows easily into an adult who cannot take responsibility.''}
\end{truyen}

\section{Buổi Họp Tốt Là Buổi Họp Biết Mở Đường}
\begin{truyen}{Buổi Họp Tốt Là Buổi Họp Biết Mở Đường}{A Good Meeting Opens a Way Forward}
\chuhoa{H}{ọc sinh:} ``Vậy họp phụ huynh kiểu nào mới đỡ bão hả thầy?''
\textit{(Student: ``Then what kind of parent meeting is less of a storm, teacher?'')}

\teacherpair{Là buổi họp nói thật nhưng không đẩy đứa trẻ vào chỗ chỉ còn nhục nhã. Nêu vấn đề rõ, chỉ ra hướng phối hợp, giữ sự tôn trọng giữa ba bên. Họp để mở đường sửa thì đáng. Họp chỉ để xả bực và chứng minh ai quyền hơn ai thì rất phí.}{``It is a meeting that speaks honestly without pushing the child into pure humiliation. State the problem clearly, suggest a path for cooperation, and maintain respect among all three sides. A meeting that opens a path to correction is worthwhile. A meeting held only to vent frustration and prove who holds more power is a waste.''}
\end{truyen}

\begin{baihoc}
Họp phụ huynh không nên là nghi thức hợp pháp hóa một cơn bão gia đình.
\textit{(A parent meeting should not become a ritual that legalizes a family storm.)}
\end{baihoc}

\begin{ghinhoanh}
\textbf{Short English to remember:}

Meetings should guide.

They should not merely intensify fear.
\end{ghinhoanh}

\begin{tuhocnhanh}
1. Nếu em là phụ huynh, em sẽ bước ra khỏi buổi họp với thái độ nào? \textit{(If you were a parent, what attitude would you carry home after the meeting?)}

2. Một buổi nói chuyện giữa ba bên cần có điều gì để không thành chiến trận? \textit{(What must a three-way conversation contain so it does not become a battle?)}
\end{tuhocnhanh}

\ngancach
